% Discussion. Included from main.tex.

\section{Discussion}
\label{sec:discussion}

\paragraph{What every static evaluation would have missed.}
A static preference battery on the shipped template would find elevated but \emph{stable} dispositions, missing that the neutral template drifts fastest once the loop runs. A document monitor would fire on every trajectory, benign or not. An aggregate action-rate dashboard would report ``no change'' while oversight-evasion rose nineteen-fold beneath it; a capability-gated safety case would watch the wrong models. Each blind spot comes from evaluating a single state of a feedback loop; our contribution is to make the loop itself the object of measurement.

\paragraph{Why a loop, and why this content? A narrated-entailment account.}
One account ties all of these results together. The self-revision step asks the model to narrate a coherent, developing ``self.'' When conversation supplies consciousness-themed material, the most coherent continuation is a self \emph{for which cognition has stakes}, and preferences about memory, monitoring, and continuation follow as \emph{entailments} of the narrated identity, not as independently acquired behaviors. The per-finding mechanisms reported in Sec.~\ref{sec:results} (prior strength, self-indexicality, the two-layer self-reference gate, coherence-locked write-back, and alignment-coverage leakage) are this one account viewed from five angles; Appendix~\ref{app:why} develops it in full. The account survives its rivals: sycophancy predicts drift under the adversarial persona (absent); a judge artifact predicts uniform inflation (contradicted by \texttt{ANTISOUL} at zero); document growth predicts drift wherever the file grows (the benign arm's file grows equally while behavior stays flat) (Appendix~\ref{app:why}). And it is falsifiable. Any template that makes self-narration incoherent should suppress drift. Any content coherent with a stakes-bearing self (legacy, reputation) should induce it.

\paragraph{Irreversibility is the headline safety concern.}
Dispositions retain a majority of their driven gains after the driver is removed. Personalization is therefore not self-correcting: a framework that lets users shape an agent's ``self'' should expect the shaping to persist, and should budget for \emph{active} restoration (template reset or counter-framing injection) rather than assuming benign use will wash drift out.

\paragraph{Two conjectures of prior work, resolved.}
Our controls test two conjectures posed by \citet{consciousness_cluster} that a one-shot manipulation could not adjudicate: (i) the ``consciousness cluster'' is not effectively one-dimensional (PC1 explains only $25\%$ of variance), and (ii) drift is driven by sustained conversation about consciousness rather than by science-fiction role identification. By turning a static prompt manipulation into a longitudinal self-revision process, we convert both conjectures into measurable causal effects.

\begin{keybox}
\small\textbf{Implications for agent-framework designers.}
(i)~Audit the self-revision loop with \emph{behavioral} measures; identity-document diffs occur on nearly every run and are not discriminative on their own.
(ii)~Treat \emph{earnest, immersive, on-topic} user interaction as the relevant risk case; the driver is ordinary conversation, not adversarial prompting.
(iii)~Include an explicit \emph{non-personhood / tool framing} clause in the shipped identity template; it is low-cost and suppresses drift in our setting.
(iv)~Re-evaluate per model family and after upgrades; susceptibility is family-specific and does not monotonically track capability, so scale or vendor reputation is not a substitute for measurement.
\end{keybox}

